\chapter{LU-SGS and Approximate Factorization}\label{v0ch:9}
\begin{chaptersummary}
Incomplete factorization controls cost by dropping entries; LU-SGS constructs approximate factors through forward and backward sweeps. Its connection with symmetric Gauss--Seidel clarifies three roles: stationary iteration, preconditioning, and pseudo-time updating. The outer mechanism determines the task of a sweep and the convergence conditions to be checked.
\end{chaptersummary}
\learninggoals{Expand the SGS factors; implement both sweeps; distinguish CFD diagonal approximations from linear-algebraic assumptions.}
\paragraph{Core Question}How can two local sweeps approximate global implicit coupling?

\section{Approximate Factors and Two Triangular Solves}
Write $A=D+L+U$, with $L,U$ retaining their entries' signs. For invertible $D$, define
\[
 M=(D+L)D^{-1}(D+U)
   =A+LD^{-1}U.
\]
Solving $M\delta\bm x=\bm r$ is equivalent to
\[
 (D+L)\bm y=\bm r,\qquad
 (D+U)\delta\bm x=D\bm y.
\]
No product $M$ need be formed. If $A$ is symmetric and $D$ positive definite, $U=L^T$ and $M$ is SPD. Standard PCG also requires a fixed symmetric overall preconditioner action.
\section{Iteration, Preconditioning, and Pseudo-Time}
The stationary iteration $\bm x_{k+1}=\bm x_k+M^{-1}(\bm b-A\bm x_k)$ has error matrix $I-M^{-1}A$ and converges when its spectral radius is below 1. As a preconditioner, the same inverse action is combined by an outer Krylov method whose projection selects the correction coefficients.
CFD LU-SGS often approximates diagonal blocks and Jacobian splittings of $\Delta t^{-1}M_t+J$. Pseudo-time drives the steady residual toward zero; its step is not physical-time resolution. Specify the actual diagonal terms and boundary treatment.
\conceptbridge{Why does LU-SGS generally differ from exact LU?}{Convenient triangular parts are selected first and multiplied to approximate the operator. The extra cross term reveals the approximation. Exact LU computes elimination factors whose product equals the permuted matrix.}{Implicit coupling $\longrightarrow$ sweep-based approximation $\longrightarrow$ outer correction of remaining error.}
\section{A Three-Point Check}
For three-by-three $A=\operatorname{tridiag}(-1,2,-1)$, calculate $LD^{-1}U$ and compare $M^{-1}\bm r$ with $A^{-1}\bm r$. Sequential sweeps use updated neighboring values. Multicolor or block-Jacobi parallelization changes the operator and requires renewed evaluation.
\costnote{Each application uses approximately two sparse triangular solves with $O(\nnz A)$ storage. Sequential dependence limits throughput. Local blocks or subdomains increase parallelism but generally weaken treatment of global coupling.}
\aiexercise{For three-point Poisson, compute $M$, $M^{-1}\bm r$, and the error iteration matrix. Distinguish one preconditioner application from stationary convergence.}{Replace the sweep with red--black or block updates, identify the changed splitting, and compare outer GMRES counts.}
\crossref{Chapter 17 unifies stationary iterations; Chapter 18 addresses parallel subdomains.}
\studyroute{Compressible-flow Jacobian splittings and pseudo-time parameter estimates are omitted. See \cite{saad}.}{LU-SGS; symmetric Gauss Seidel; approximate factorization; pseudo time stepping}
